% Problem record op_7920f48995bc8511. This TeX file is derived from
% database/problems_json/op_7920f48995bc8511.json by scripts/sync-tex.mjs; edit the JSON record
% and rerun the script rather than editing this file.
\section{Nontrivial mutually degradable channel pairs}
\label{sec:op_7920f48995bc8511}

\paragraph{Problem.}
Does there exist an integer $d\geq2$ and a pair of distinct channels
$\mathcal M,\mathcal N:\mathcal L(A)\to\mathcal L(B)$, with
$A\simeq B\simeq\mathbb C^d$, that both have Choi rank exactly $d$, are
mutually degradable, and are each nondegradable?  Let
$E\simeq\mathbb C^d$ and choose minimal Stinespring isometries
$V_{\mathcal M},V_{\mathcal N}:A\to B\otimes E$ defining the channels and
their complements by
\begin{equation}
  \begin{aligned}
    \mathcal M(\rho)&=\operatorname{Tr}_E
      (V_{\mathcal M}\rho V_{\mathcal M}^{\dagger}),
    &\mathcal M^c(\rho)&=\operatorname{Tr}_B
      (V_{\mathcal M}\rho V_{\mathcal M}^{\dagger}),\\
    \mathcal N(\rho)&=\operatorname{Tr}_E
      (V_{\mathcal N}\rho V_{\mathcal N}^{\dagger}),
    &\mathcal N^c(\rho)&=\operatorname{Tr}_B
      (V_{\mathcal N}\rho V_{\mathcal N}^{\dagger}).
  \end{aligned}
  \label{eq:p59-complementary-pairs}
\end{equation}
Equation~\eqref{eq:p59-complementary-pairs} fixes representatives of the
complementary channels; changing a minimal dilation only applies an output
unitary to a complement.

For
$\lvert\Omega_d\rangle:=\sum_{j=1}^d\lvert j\rangle_{A'}\lvert j\rangle_A$,
the required Choi-rank condition is
\begin{equation}
  J(\mathcal T):=(\operatorname{id}_{A'}\otimes\mathcal T)
    (\lvert\Omega_d\rangle\!\langle\Omega_d\rvert),
  \qquad
  \operatorname{rank}J(\mathcal M)
   =\operatorname{rank}J(\mathcal N)=d.
  \label{eq:p59-choi-rank}
\end{equation}
The equality in Eq.~\eqref{eq:p59-choi-rank} makes the environment dimension
in Eq.~\eqref{eq:p59-complementary-pairs} minimal.

Mutual degradability requires channels
$\mathcal X,\mathcal Y:\mathcal L(B)\to\mathcal L(E)$ such that
\begin{equation}
  \mathcal X\circ\mathcal M=\mathcal N^c,
  \qquad
  \mathcal Y\circ\mathcal N=\mathcal M^c.
  \label{eq:p59-mutual-degradability}
\end{equation}
In addition to Eq.~\eqref{eq:p59-mutual-degradability}, neither channel may
admit its own degrading map:
\begin{equation}
  \begin{aligned}
    &\nexists\ \mathcal D_{\mathcal M}:\mathcal L(B)\to\mathcal L(E)
      \quad\text{CPTP with}\quad
      \mathcal M^c=\mathcal D_{\mathcal M}\circ\mathcal M,\\
    &\nexists\ \mathcal D_{\mathcal N}:\mathcal L(B)\to\mathcal L(E)
      \quad\text{CPTP with}\quad
      \mathcal N^c=\mathcal D_{\mathcal N}\circ\mathcal N.
  \end{aligned}
  \label{eq:p59-individual-nondegradability}
\end{equation}
Equation~\eqref{eq:p59-individual-nondegradability}, together with
$\mathcal M\neq\mathcal N$, excludes the identity-channel and self-pair constructions described below.
It does not exclude complementary pairs.

\subsection*{Status}

Solved

\subsection*{Source}

Ruskai posed mutual degradability in Problem 23, Eq.~(32), and singled out
two Choi-rank-$d$ channels that are not individually degradable
\sourcecite{ref:p59-ruskai}{Rus07}. The existence statement is answered by
the explicit complementary qutrit pair recorded in item A of the catalog's
scientific-review issue \sourcecite{ref:p59-review-41}{Rev26}.

\subsection*{Progress}

\begin{itemize}
  \item An explicit solution has $d=3$, $a=3/4$, and $b=1/3$. Define the
isometry
\begin{equation}
  \begin{aligned}
    V|0\rangle&=|00\rangle,\\
    V|1\rangle&=\sqrt{1-a}|10\rangle+\sqrt a|01\rangle,\\
    V|2\rangle&=\sqrt{1-b}|20\rangle+\sqrt b|02\rangle.
  \end{aligned}
  \label{eq:p59-qutrit-isometry}
\end{equation}
Set $\mathcal M=\operatorname{Tr}_E V(\cdot)V^\dagger$ and
$\mathcal N=\operatorname{Tr}_B V(\cdot)V^\dagger$ in
Eq.~\eqref{eq:p59-qutrit-isometry}, identifying both output spaces with
$\mathbb C^3$. Use the swapped isometry for $\mathcal N$. Then
$\mathcal M^c=\mathcal N$ and $\mathcal N^c=\mathcal M$, so
$\mathcal X=\mathcal Y=\operatorname{id}$ satisfies
Eq.~\eqref{eq:p59-mutual-degradability}.
The Kraus operators of $\mathcal M$ are
$\operatorname{diag}(1,\sqrt{1-a},\sqrt{1-b})$,
$\sqrt a|0\rangle\langle1|$, and $\sqrt b|0\rangle\langle2|$;
those of $\mathcal N$ replace $(a,b)$ by $(1-a,1-b)$.
Both triples are linearly independent, proving
Eq.~\eqref{eq:p59-choi-rank}, and the outputs on
$|1\rangle\langle1|$ differ.
If degrading maps $\mathcal D\mathcal M=\mathcal N$ and
$\mathcal D'\mathcal N=\mathcal M$ existed, their actions on the ground
state and the relevant excited state would force
\begin{equation}
  \begin{aligned}
    \mathcal D(|1\rangle\langle1|)
      &=3|1\rangle\langle1|-2|0\rangle\langle0|,\\
    \mathcal D'(|2\rangle\langle2|)
      &=2|2\rangle\langle2|-|0\rangle\langle0|.
  \end{aligned}
  \label{eq:p59-nonpositive-degraders}
\end{equation}
Neither output in Eq.~\eqref{eq:p59-nonpositive-degraders} is positive,
proving Eq.~\eqref{eq:p59-individual-nondegradability}. This is the
unpublished construction from the scientific-review issue
\sourcecite{ref:p59-review-41}{Rev26}.

  \item At unrestricted rank, Ruskai observed that the identity channel and an
  arbitrary channel form a mutually degradable pair.  Also, any degradable
  channel paired with itself satisfies
  Eq.~\eqref{eq:p59-mutual-degradability}.  The rank, distinctness, and
  nondegradability requirements exclude both constructions
  \sourcecite{ref:p59-ruskai}{Rus07}.

  \item Cubitt, Ruskai, and Smith proved that every qubit channel with two
  Kraus operators is either degradable or antidegradable.  Consequently, any
  $d=2$ solution satisfying Eq.~\eqref{eq:p59-individual-nondegradability}
  must consist of two antidegradable channels.  Their classification neither
  constructs nor excludes such a pair satisfying
  Eq.~\eqref{eq:p59-mutual-degradability}
  \sourcecite{ref:p59-cubitt-ruskai-smith}{CRS08}.
\end{itemize}

\subsection*{References}

\begin{enumerate}[label={},leftmargin=4.8em,labelsep=0.5em]
  \footnotesize
  \item[\textup{[Rus07]}]\label{ref:p59-ruskai}
  M. B. Ruskai, ``Open Problems in Quantum Information Theory,''
  arXiv preprint arXiv:0708.1902 (2007).
  \newline
  \href{https://doi.org/10.48550/arXiv.0708.1902}{doi:10.48550/arXiv.0708.1902};
  \href{https://arxiv.org/abs/0708.1902}{arXiv:0708.1902}.

  \item[\textup{[CRS08]}]\label{ref:p59-cubitt-ruskai-smith}
  T. S. Cubitt, M. B. Ruskai, and G. Smith,
  ``The Structure of Degradable Quantum Channels,''
  \emph{Journal of Mathematical Physics} \textbf{49}, 102104 (2008).
  \href{https://doi.org/10.1063/1.2953685}{doi:10.1063/1.2953685};
  \href{https://arxiv.org/abs/0802.1360}{arXiv:0802.1360}.

  \item[\textup{[Rev26]}]\label{ref:p59-review-41}
  QIQCOP Zoo, “Scientific review: confirm channel, resource and QEC
statements in 10 catalog records,” GitHub issue \#41, item A (2026).
Unpublished construction.
\href{https://github.com/Naixu-Guo/quantum-open-problems/issues/41}{Issue \#41}.
\end{enumerate}

\subsection*{Comment}

The displayed existence question is solved by
Eqs.~\eqref{eq:p59-qutrit-isometry} and
\eqref{eq:p59-nonpositive-degraders}. The resolving construction is an
unpublished issue contribution, not a peer-reviewed result. Excluding
complementary pairs would define a stronger question; no such exclusion
appears in the archived statement or in Ruskai's Problem 23.

\subsection*{Field}

Quantum Communication

\subsection*{Topic}

Channel degradability; Quantum channel structure

\subsection*{ID}

\texttt{op\_7920f48995bc8511}
